\chapter{System Implementation}

\section{Development Environment}
\begin{sloppypar}
\raggedright
The proposed implementation uses a Manifest V3 extension for Chromium-based browsers,
React with Tailwind CSS for the dashboard, Python with FastAPI for backend services,
Firebase Authentication for identity. Python
machine-learning components use XGBoost for URL and file classification and Random
Forest for initial email screening. QR-decoding, file-analysis, Playwright testing,
SHAP explanations, and configured RAG/LLM or threat-intelligence services are added
where required by the selected prototype modules.
\end{sloppypar}

\section{Component Integration}
The extension communicates with FastAPI through HTTPS and sends Firebase ID tokens
for protected requests. FastAPI verifies the token, checks authorization, validates
the scan request, invokes the appropriate detector, and returns a common result shape.
The dashboard uses the same API rather than connecting directly to detector services. Detector
modules remain independently testable so that models, rules, integrations, and
retention policies can be updated without changing the user interface.

\section{Detector Implementation}
URL detection normalizes inputs and extracts configured lexical and domain features.
Page detectors use content scripts and supported browser APIs to inspect scripts,
redirects, frames, notification prompts, and observable DOM behaviour. Download
analysis uses content-based file identification, format-specific parsers, hash and
reputation checks, resource limits, and isolated processing. Email screening is
lightweight in the extension and uses RAG/LLM analysis only for an authorized detailed
dashboard investigation. Playwright is reserved for controlled analysis and automated
testing rather than indiscriminate collection of user credentials.

\section{User Interface and Reports}
The extension presents the risk level, reason, evidence, and available action without
overwhelming non-technical users. The dashboard provides login, URL/QR/email/file scan
forms, detection history, threat reports, analytics, and troubleshooting information.
Reports distinguish detection from prevention and identify incomplete analysis,
unsupported formats, timeouts, and unavailable browser capabilities.

\section{Testing and Evaluation}
Unit tests will validate feature extraction, parsers, rules, risk decisions, token
verification, authorization, and privacy filtering. Integration tests will cover
extension-to-API communication, authorization checks, model invocation, and
report generation. Browser tests will use safe synthetic pages, controlled phishing
fixtures, and synthetic credentials. Evaluation will report precision, recall,
false-positive and false-negative rates, latency, resource consumption, model
generalization, and confirmed versus requested protective actions.

\section{Deployment and Operations}
The extension and dashboard require configured HTTPS endpoints, Firebase settings,
backend access restricted to server-side services, model artifacts, and external-service
credentials stored outside client code. Logging will contain detection evidence and
actions without unnecessary sensitive content. Configuration will define file-size and
archive limits, timeouts, retention, supported browsers, and the policy for warning,
cancellation, quarantine, or release.
